% ==============================================================================
% CAPÍTULO 3 - REFERENCIAL TEÓRICO
% ==============================================================================

\section{Referencial teórico}
\label{sec:referencial-teorico}

Esta seção reúne os conceitos em que o trabalho se apoia, organizados em dois eixos: os
fundamentos (segurança de bancos de dados, o registro de log do PostgreSQL, detecção de
anomalias com o \iforest{} e os métodos de avaliação) e os trabalhos relacionados à detecção de
acessos anômalos em bancos de dados relacionais, agrupados pela referência de normalidade que
adotam.

\subsection{Segurança de bancos de dados e o log do PostgreSQL}
\label{sec:seguranca-log}

O principal mecanismo de confidencialidade de um banco de dados é o controle de acesso, que
decide, a cada operação, se o sujeito tem permissão sobre o objeto. \citeonline{bertino2005}
descrevem os modelos clássicos: discricionário, mandatório e baseado em papéis (RBAC). O
controle de acesso responde a uma pergunta binária e não avalia o contexto em que uma operação
permitida acontece. Esse é o terreno das ameaças internas: nos mais de setecentos casos reunidos
por \citeonline{cappelli2012}, o que denuncia o problema é o uso que se faz do acesso, e não a
falta de permissão. A detecção de comportamento anômalo complementa o controle de acesso ao
perguntar se a operação é típica para quem a executa.

O PostgreSQL registra eventos do servidor de acordo com parâmetros de configuração
\cite{postgresql2025a}. Com \texttt{log\_destination = jsonlog}, cada entrada vira um objeto
JSON com horário, usuário, banco, aplicação, endereço de origem, sessão, severidade, mensagem e
texto do comando; com \texttt{log\_min\_duration\_statement = 0}, todo comando concluído é
registrado; \texttt{log\_connections} e \texttt{log\_disconnections} registram conexões, e as
falhas de autenticação são sempre registradas. O log nativo, porém, não informa quantas linhas
um comando retornou ou alterou. Essa informação vem do módulo \texttt{auto\_explain}, que
registra automaticamente os planos de execução \cite{postgresql2025b}: com
\texttt{log\_analyze}, o plano traz a contagem real de linhas de cada nó e o nome das tabelas
acessadas.

\subsection{Detecção de anomalias e o \textit{Isolation Forest}}
\label{sec:anomalias-iforest}

\citeonline{chandola2009} definem anomalias como padrões que não estão de acordo com uma noção
bem definida de comportamento normal, e distinguem anomalias pontuais, contextuais e coletivas.
Um acesso às três da manhã é normal para uma rotina noturna e atípico para um atendente de
balcão: é uma anomalia contextual, tratada neste trabalho ao comparar cada usuário com o próprio
histórico. Como rótulos de acessos indevidos são escassos e cobrem só ataques já conhecidos, a
abordagem não supervisionada é a mais natural. Ela enfrenta, porém, o desbalanceamento entre
classes, a deriva do comportamento normal ao longo do tempo \cite{gama2014} e a distância entre
o laboratório e o uso real \cite{sommer2010}.

O \iforest{} parte da ideia de que anomalias são ``poucas e diferentes'' e, por isso, mais fáceis
de isolar \cite{liu2008,liu2012}. Cada árvore é construída sobre uma subamostra de $\psi$
observações: em cada nó, sorteiam-se um atributo e um valor de corte entre o mínimo e o máximo
desse atributo, até isolar cada observação. Uma observação atípica tende a ser isolada perto da
raiz, com um caminho curto. O comprimento médio do caminho é normalizado pelo de uma busca
malsucedida em uma árvore binária com $n$ elementos,

\begin{equation}
    c(n) = 2H(n-1) - \frac{2(n-1)}{n},
    \label{eq:cn}
\end{equation}

\noindent em que $H(i)$ é o número harmônico, aproximado por $\ln(i) + 0{,}5772156649$. O escore
de anomalia de uma observação $x$ é

\begin{equation}
    s(x, n) = 2^{-\frac{E(h(x))}{c(n)}},
    \label{eq:score}
\end{equation}

\noindent em que $E(h(x))$ é a média do comprimento do caminho de $x$ nas árvores. Escores
próximos de 1 indicam anomalias; escores bem abaixo de 0,5, observações normais. O algoritmo
não usa distância nem densidade, tem custo proporcional a $t\,\psi \log \psi$ para $t$ árvores e
funciona bem com subamostras pequenas \cite{liu2012}. A implementação usada é a da biblioteca
scikit-learn \cite{pedregosa2011}, em que o método \texttt{score\_samples} devolve o oposto do
escore original e o parâmetro \texttt{contamination} só define o limiar de decisão, sem alterar
as árvores \cite{scikitlearn2026}. Duas limitações aparecem nos resultados: os cortes paralelos
aos eixos produzem artefatos em regiões sem dados \cite{hariri2021}, e, como os cortes ficam entre
o mínimo e o máximo do treino, o escore satura para valores muito acima do observado.

\subsection{Avaliação de detectores}
\label{sec:avaliacao-detectores}

A partir da matriz de confusão, a precisão é a fração dos alertas que eram anomalias, o recall é
a fração das anomalias que geraram alerta e o F1 é a média harmônica dos dois. Quando o detector
produz um escore, as áreas sob as curvas ROC e precisão--recall (PR-AUC) o avaliam sem fixar um
limiar; com classes desbalanceadas, a curva precisão--recall é mais informativa
\cite{saito2015}.

Para comparar detectores, não basta a diferença entre médias. O \textit{bootstrap} percentil
estima intervalos de confiança reamostrando os dados \cite{efron1993}; quando as observações são
dependentes, como janelas consecutivas de um mesmo usuário, reamostram-se blocos de observações
\cite{kunsch1989}. O teste de McNemar compara dois classificadores nos mesmos casos, usando só os
discordantes, e tem baixa taxa de erro do tipo~I \cite{dietterich1998}. O teste de permutação
pareado constrói a distribuição nula trocando as saídas dos dois detectores dentro de cada par
\cite{noreen1989}, o que serve a métricas como o F1. Com várias hipóteses, a correção de
\citeonline{holm1979} controla a taxa de erro da família. Por fim, o pré-registro separa a
geração de hipóteses do seu teste e evita que a pergunta mude depois de vista a resposta
\cite{nosek2018}.

\subsection{Trabalhos relacionados}
\label{sec:trabalhos-relacionados}

\textbf{Perfis por papel a partir da sintaxe.} \citeonline{kamra2008} mineram as consultas SQL
do log de auditoria para construir perfis de acesso normal: com controle de acesso por papéis,
um classificador aponta o usuário cujas consultas não se parecem com as do seu papel; sem
papéis, as consultas são agrupadas e os grupos fazem o papel dos perfis.
\citeonline{shebaro2013} levam essa ideia para dentro do PostgreSQL, com atributos da árvore
sintática de cada comando e um classificador bayesiano que estima o papel mais provável; se ele
difere do papel de quem executou o comando, o sistema emite um alarme.
\citeonline{ronao2016} seguem a mesma formulação, com floresta aleatória e análise de
componentes principais. Essas abordagens perguntam se a consulta se parece com as do papel e não
captam, por exemplo, um atendente que faz consultas típicas de atendente, mas em volume, horário
ou origem fora do seu costume.

\textbf{Perfis a partir dos dados acessados.} \citeonline{mathew2010} partem da premissa de que a
sintaxe da consulta, sozinha, é um sinal fraco da intenção do usuário, e modelam o padrão de
acesso de cada usuário pelos dados que as suas consultas efetivamente acessam. Essa preocupação
com o que foi lido também está no SACLA, que usa as linhas retornadas e as tabelas acessadas.

\textbf{Sistemas de detecção e resposta.} \citeonline{sallam2017} apresentam o DBSAFE, que
constrói perfis do comportamento normal de usuários e aplicações, detecta acessos anômalos e
responde a eles com políticas predefinidas, automáticas ou humanas. A resposta automática está
fora do escopo deste trabalho, em que a decisão cabe sempre a uma pessoa. Na avaliação de
detectores, \citeonline{goldstein2016} compararam 19 algoritmos não supervisionados em 10
conjuntos de dados, e \citeonline{sommer2010} recomendam que a avaliação explique o que o
detector encontra e reconheça as limitações dos dados usados.

O \autoref{qdr:trabalhos} resume a posição do SACLA. Ele se diferencia por usar só o log nativo,
sem modificar o SGBD; por agregar a atividade de cada usuário em janelas e compará-la com o
histórico do próprio usuário, o que capta volume, horário, origem, erros e falhas de
autenticação; por explicar cada alerta; e por uma avaliação com linha de base, réplicas e
hipóteses pré-registradas.

\begin{quadro}[htb]
    \centering
    \ABNTEXfontereduzida
    \begin{minipage}{\textwidth}
        \caption{\label{qdr:trabalhos}Comparação entre este trabalho e os trabalhos relacionados}
        \begin{tabular}{|L{2.8cm}|L{3.2cm}|L{2.0cm}|L{3.0cm}|L{2.6cm}|}
            \hline
            \textbf{Trabalho} & \textbf{Fonte de dados} & \textbf{Unidade} & \textbf{Aprendizado} & \textbf{Normalidade} \\
            \hline
            \citeonline{kamra2008} & log de auditoria das consultas & consulta & classificação ou agrupamento & papel ou grupo \\
            \hline
            \citeonline{mathew2010} & dados acessados pelas consultas & consulta & aprendizado estatístico & perfil de acesso aos dados \\
            \hline
            \citeonline{shebaro2013} & árvore sintática, dentro do PostgreSQL & consulta & classificador bayesiano & papel \\
            \hline
            \citeonline{ronao2016} & consultas em banco com papéis & consulta & floresta aleatória e PCA & papel \\
            \hline
            \citeonline{sallam2017} & perfis de usuários e aplicações & consulta & perfis aprendidos em treino & usuário ou aplicação \\
            \hline
            Este trabalho & log nativo (\textit{jsonlog} e \texttt{auto\_explain}) & janela de 5~min por usuário & \iforest{}, não supervisionado & histórico do próprio usuário \\
            \hline
        \end{tabular}
        \legend{Fonte: Própria do autor (2026).}
    \end{minipage}
\end{quadro}
